\begin{longtable}{@{}p{1.5cm}p{4.6cm}p{1.1cm}p{1.5cm}p{5.0cm}@{}}
\caption{Korpus basis pengetahuan.}
\label{tab:korpus-kb} \\
\toprule
\textbf{Kode} & \textbf{Lembaga} & \textbf{Tahun} & \textbf{Halaman} & \textbf{Topik utama} \\
\midrule
\endfirsthead

\multicolumn{5}{@{}l}{\small\emph{Tabel \thetable{} (lanjutan)}} \\
\toprule
\textbf{Kode} & \textbf{Lembaga} & \textbf{Tahun} & \textbf{Halaman} & \textbf{Topik utama} \\
\midrule
\endhead

\bottomrule
\endlastfoot

KB-01 & IDAI UKK Endokrinologi Anak dan Remaja & 2015 & 129 & Hipoglikemia, insulin \\
KB-02 & PERKENI & 2021 & 49 & Pemantauan, insulin \\
KB-03 & PERKENI & 2021 & 70 & Insulin \\
KB-04 & PERKENI & 2023 & 105 & Hiperglikemia, insulin \\
KB-05 & ADA dan EASD & 2021 & 44 & Hipoglikemia, insulin, target glikemik \\
KB-06 & Konsensus Internasional ATTD & 2019 & 11 & Pemantauan, target glikemik \\
KB-07 & ISPAD Bab 10 & 2022 & 25 & Karbohidrat, insulin \\
KB-08 & ISPAD Bab 11 & 2022 & 22 & Ketoasidosis, insulin \\
KB-09 & ISPAD Bab 12 & 2022 & 19 & Hipoglikemia, insulin \\
KB-10 & ISPAD Bab 13 & 2022 & 14 & Ketoasidosis, insulin, hari sakit \\
KB-11 & ISPAD Bab 14 & 2022 & 32 & Aktivitas, karbohidrat, hipoglikemia \\
KB-12 & ISPAD Bab 25 & 2022 & 23 & Insulin, penyesuaian tata laksana pada keterbatasan sumber daya, bersifat lintas kondisi \\
\midrule
& \textbf{Total} & & \textbf{543} & \\

\end{longtable}

\noindent\footnotesize Judul lengkap setiap dokumen disajikan pada Lampiran. Kolom topik utama disarikan dari manifes korpus, yang memuat pula nilai penyeimbang nomor halaman untuk keperluan keterlacakan sitasi sebagaimana dituntut KNF-08.\normalsize
